\subsubsection*{A determinant supported on the theta points}
\label{subsec:congruence-determinant}
\hypertarget{source-congruence-determinant}{}

The projections and inertia relations describe the relevant blocks of the integral Galois matrix algebra. A determinant in the Hecke algebra will convert their cohomology lengths into a congruence estimate supported on the theta points. Fix one coefficient DVR \(O\) with
fraction field \(F\). In this calculation write
\(\operatorname{len}\) for the unnormalized \(O\)-length.
Use a finite parameter set \(X\), the theta subset \(Z\), a
Hecke algebra \(\mathcal T\subset O^X\) and a good point
\(\pi\in Z\). Write \(I_Z\) for the ideal of elements of
\(\mathcal T\) vanishing outside \(Z\). The congruence
length before normalization is
\(\operatorname{len}(O/I_Z(\pi))\).

\begin{established}[Ordered corner and block evaluation]\label{hyp:theta-marked-corner}
The auxiliary parameter set \(X'\) records an ordering of
the eigenvalues of a marked \(t\)-dimensional corner. Its
coefficient algebra \(\mathcal T'\subset O^{X'}\) has a
basis \(b_j\) over the pullback of \(\mathcal T\) along
\(p:X'\to X\). For each \(x\in X\setminus Z\), evaluation gives an injection
\[
 O^{t!}\lhook\joinrel\longrightarrow O^{p^{-1}(x)},\qquad
 (c_j)_j\longmapsto\bigg(\sum_j c_j b_j(x')\bigg)_{x'\in p^{-1}(x)} .
\] At the chosen
point of \(X'\) above \(\pi\) the matrix algebra
\(\mathcal A\) has a specialization into matrices over
\(O\). Its image is block triangular with one block of
size \(r\) and \(t\) blocks of size one, and each block
map \(\rho_k:\mathcal A\to M_{d_k}(O)\) is surjective.
The determinants of square matrices over \(\mathcal A\),
computed on the full matrix spaces at all points of
\(X'\), belong to \(\mathcal T'\).
\end{established}
\begin{proof}
\proofstep{Integral coefficients of the matrix determinants.}
These integral properties can be seen from the marked
corner construction. Let \(\mathcal A_0\) be the
\(\mathcal T\)-span of the Galois matrices \(r(g)\),
which is an algebra since products of Galois matrices
are Galois matrices. The coefficients of the
characteristic polynomials of \(r(g)\) are in
\(\mathcal T\). For any finite number of block rows,
expand a matrix over \(\mathcal A_0\) as a scalar linear
combination of \(e_{ij}\otimes r(g)\), with \(e_{ij}\)
the elementary matrices on the block rows. To see that
its characteristic coefficients are in \(\mathcal T\),
give each summand \(A_i=e_{ij}\otimes r(g)\) an
indeterminate \(X_i\). For a primitive cyclic word \(w\), let
\(A_w\) be its ordered matrix product and \(X_w\) its monomial.

In the completed polynomial ring in characteristic zero,
\[
 \log\det(1-\textstyle\sum_iX_iA_i)
 =-\sum_{a\ge1}\frac{\operatorname{tr}((\sum_iX_iA_i)^a)}a
 =-\sum_{w\ {\rm primitive}}\sum_{b\ge1}
                    \frac{X_w^b\operatorname{tr}(A_w^b)}b .
\]
Indeed \(w^b\) contributes \(|w|\) rotations at length \(b|w|\).
Exponentiating gives \(\prod_w\det(1-X_w A_w)\), with only
finitely many factors in each degree.
A product \(A_w\) is zero, has one off-diagonal
block, or has a single diagonal block carrying a
Galois matrix. The individual factors are therefore
\[
 \det(1-X_w A_w)\in\mathcal T[X_w].
\]
For the zero and off-diagonal cases the determinant equals one. After the identity has been
proved in the fraction fields, its coefficient
expression is integral and uses no division in
\(\mathcal T\). Substitution of the scalar coefficients
proves the assertion, also over any \(\mathcal T\)-algebra
of scalar functions.

\proofstep{The ordered corner roots and their splitting algebra.}
The corner operator in \(\mathcal A_0\) is
\(e_Br(g_0)e_B+s\), where \(e_B\) is the marked
rank-\(t\) idempotent and \(s=e_Bse_B\) is the
perturbation with \(s(\pi)=0\). It is zero on the
complement of the corner. The previous argument
places its corner characteristic polynomial in
\(\mathcal T[X]\), with \(X\) here a polynomial
indeterminate. The marked perturbation gives distinct
roots at every fibre in the tied splitting field.
The successive monic quotients adjoining the ordered roots have degrees \(t,t-1,\ldots,1\). Freeness at each step gives
\[
 \operatorname{rank}_{\mathcal T}(\text{splitting algebra})=\prod_{a=1}^{t}a=t!.
\]

Over each rational fibre this algebra is the
product of \(t!\) copies of the coefficient field,
by the Chinese remainder theorem. The free basis
therefore stays independent when evaluated at all
orderings in each fibre. This also gives an embedding
of the splitting algebra into \(O^{X'}\), since the
roots of integral monic polynomials are integral.
Take \(\mathcal A\) to be the span of the Galois
matrices over this splitting algebra. The determinant
argument applies with these coefficients. All
matrices are integral, so \(\mathcal A\) is finite
over \(O\). At the chosen fibre the marked
decomposition gives the triangular blocks, and
the residual full matrix span of the rank-\(r\)
constituent gives surjectivity on that block by
Nakayama's lemma. The scalar blocks have surjectivity
from the coefficient ring. These algebraic steps
use the marked idempotents and the perturbation
on the given family, constructed in Proposition~\ref{prop:retained-integral-projections}.
\end{proof}

\begin{proposition}[The congruence ideal on the theta locus]
\label{prop:theta-determinant-length}
Take the ordered corner and block maps of Lemma~\ref{hyp:theta-marked-corner}. Let \(M\) be the finite
presentation with the column relations in the proof. If these
relations have the prescribed finite column dimensions away from
\(Z\), then
\[
 \operatorname{len}(O/I_Z(\pi))\le
       \sum_k\operatorname{len}\operatorname{Hom}_{\mathcal A}
                          (M,(F/O)^{d_k}).
\]
The \(k\)-th action is via \(\rho_k\), with the same coefficient
ring and chosen point in all the summands.
\end{proposition}

\begin{proof}
Lemma~\ref{hyp:theta-marked-corner} supplies the splitting algebra, its fibre basis and the full block evaluations.

Write
\[
 B=e_Br(g_0)e_B+s
\]
on the corner, and let its ordered
roots be \(\alpha_0,\ldots,\alpha_{t-1}\). The operator separating
the first eigenline is
\[
 E=\prod_{i>0}(B-\alpha_i e_B),\qquad
 E|_{V_{\alpha_i}}=
 \begin{cases}\Delta\,{\rm id},&i=0,\\0,&i>0,\end{cases}
 \qquad\Delta=\prod_{i>0}(\alpha_0-\alpha_i).
\]
Distinctness of the roots gives \(E^2=\Delta E\) on every fibre.
At \(\pi\) the roots are \(\chi(g_0)\varepsilon(g_0)^{-i}\), where
\(\varepsilon\) is the cyclotomic character. The
marked value is \(\varepsilon(g_0)=1+\ell^s v\)
with \(v\in\mathbb Z_\ell^\times\), at \(s\ge2\) for
\(\ell=2\) and \(s\ge1\) for \(\ell=3\).
The binomial calculation gives
\[
     v_\ell(\Delta_\pi)=(t-1)s+v_\ell((t-1)!).
\]
\label{cor:congruence-determinant-1}

\proofstep{The column relations in a finite presentation.}
Here is the module used in the proof. Write \(U_g\) for the
Galois operator after dividing by the first character.
Use \(E,E_1\in\mathcal A\) with \(E^2=\Delta E\) and
\(E_1^2=E_1\), with \(\Delta\in\mathcal T'\) central.
The two operators have rank one over \(F\) at each
point of \(X'\) above \(X\setminus Z\), and
\(\Delta\) is then nonzero. Put
\(Y=\sum_g\mathcal A (U_g-1)E\subset\mathcal A\).

At each column place \(u\) with decomposition subgroup
\(D_u\), adjoin a generator \(n_u\), \(j_u\) generators
\(s_{u,a}\) and a generator \(s_{u,0}'\) to \(Y\).
Impose the relations
\[
 \begin{gathered}
 E n_u=0,\qquad (E_1-1)s_{u,a}=0,\qquad
 (E-\Delta)s_{u,0}'=0,\\
 (U_g-1)E+(U_g-1)n_u
     =\Delta\left(\sum_{a=1}^{j_u}K_{u,a}(g)s_{u,a}
                  +K_{u,0}(g)s_{u,0}'\right)\qquad(g\in D_u).
 \end{gathered}
\]
Here the first term in the last relation is the
element of \(Y\). The functions \(K_{u,a},K_{u,0}\)
are the matrix cocycles in the column data. Denote the
quotient by \(M\). Since \(\mathcal A\) is left
Noetherian, \(Y\) is finitely generated and \(M\)
has a finite presentation.

\proofstep{A kernel vector on each complementary fibre.}
The candidate determinant has integral coefficients. To lie in the theta congruence ideal it must vanish when evaluated on a complementary fibre, which can be tested by producing a vector in its kernel there.

Consider a rational evaluation space \(V\) at \(x'\)
over \(X\setminus Z\). Fix a nonzero vector \(\eta=Ew\)
in the image of \(E\) and a nonzero \(s\) in the image
of \(E_1\). The two conditions on a vector in the first image are incompatible:
\[
 \operatorname{im}E\ \cap\ \bigcap_g\ker(U_g-1)=0.
\]
At \(u\), the
column condition says that the functions
\(K_{u,a}(\cdot)s\) and \(K_{u,0}(\cdot)\eta\)
lie in a space containing the coboundaries and of
dimension at most \(j_u+\dim V^{D_u}\) modulo
coboundaries. Here invariants and coboundaries use \(U\).
If \(V^{D_u}\ne0\), take \(v_u\ne0\) invariant there
and all scalars zero. Otherwise \(V^{D_u}=0\). Let \(L_u\) be the span of the classes of the \(j_u+1\) functions in the quotient by coboundaries. The column bound gives
\[
 \dim L_u\le j_u<j_u+1.
\]

A nontrivial linear dependence in this quotient gives a nonzero tuple \((v_u,(a_{u,a})_a,a_{u,0})\)
such that
\[
 (U_g-1)v_u
      =\sum_a a_{u,a}K_{u,a}(g)s+a_{u,0}K_{u,0}(g)\eta
           \qquad(g\in D_u).
\]
If at some \(u\) we can take \(Ev_u=0\), map \(Y\)
and all other places to zero and map the generators
at \(u\) to \(\Delta v_u,a_{u,a}s,a_{u,0}\eta\).
This satisfies the relations and gives a nonzero map
\(M\to V\). In the other case rescale each tuple so
that \(Ev_u=\eta\). Map \(y\in Y\) to \(yw\) and
the local generators to
\(\Delta v_u-\eta,a_{u,a}s,a_{u,0}\eta\).
The relations follow from \(E\eta=\Delta\eta\)
and the equation at \(D_u\). The map is nonzero
on \(Y\) because \(\eta\) is moved by some \(U_g\).
Thus
\[
 \operatorname{Hom}_{\mathcal A}(M,V)\ne0
\]
at every evaluation point over \(X\setminus Z\).

\proofstep{Blockwise Smith lengths at the chosen theta point.}
Present \(M\) with \(g\) generators and a finite
matrix of relations \(P\), written so that applying
\(\operatorname{Hom}_{\mathcal A}(-,V)\) identifies it
with the kernel of \(P\) acting on \(V^g\). At the
chosen point above \(\pi\), use the block targets
\[
 H_k=(F/O)^{d_k}
\]
with action via \(\rho_k\).
If their Hom lengths are finite, the expanded relation
matrix \(P_k\) in each block has full column rank.
Let \(\varpi^{a_{k,i}}\) be its elementary divisors, with
\(1\le i\le gd_k\). Since
\[
 (F/O)[\varpi^a]\simeq O/(\varpi^a),
\]
the kernel on the divisible target and its determinantal ideal satisfy
\[
 \begin{aligned}
 \operatorname{len}\operatorname{Hom}_{\mathcal A}(M,H_k)
     &=\sum_{i=1}^{gd_k}a_{k,i}=:s_k,\\
 I_{gd_k}(P_k)&=(\varpi^{s_k}),
 \end{aligned}
\]
where \(I_{gd_k}\) is generated by the maximal minors. Thus one such
minor attains valuation \(s_k\).

For each \(k\), there exists a contraction of the rows
giving a square \(gd_k\)-matrix with determinant
valuation \(s_k\). Surjectivity of \(\rho_k\) lifts
such a contraction to a rectangular matrix \(C_k\)
over \(\mathcal A\). Use the matrix
\(C(\mathbf u)=\sum_j u_j C_j\) over a polynomial
ring with indeterminates \(u_j\).
In block \(k\), the determinant of \(CP\), divided
by a uniformizer to the power \(s_k\), is a polynomial
over \(O\). Indeed Cauchy--Binet expresses the numerator as
\[
 \det(C(\mathbf u)_kP_k)=
   \sum_{|J|=gd_k}\det(C(\mathbf u)_{k,*,J})\det(P_{k,J,*}).
\]
Each term is divisible by \(\varpi^{s_k}\) in \(O[\mathbf u]\).
Its reduction is nonzero since setting \(u_k=1\)
and the other \(u_j=0\) gives a unit. Therefore the
product of the reductions is a nonzero polynomial,
and some coefficient is a unit.

\proofstep{A determinant coefficient in the original congruence ideal.}
The same characteristic coefficient argument applied
over \(\mathcal T'[\mathbf u]\) puts the determinant of
\(C(\mathbf u)P\) in \(\mathcal T'[\mathbf u]\). At every
point above \(X\setminus Z\) the nonzero maps \(M\to V\)
give a kernel vector, so the determinant vanishes
as a polynomial at those points. At the data at
\(\pi\) the determinant is the product of the block
determinants by triangularity. Take a coefficient
\(f'\in\mathcal T'\) corresponding to a unit coefficient
of the reduced product. Its valuation at this point
is \(\sum_k s_k\), and it vanishes off the theta subset.
Write
\[
 f'=\sum_j(f_j\circ p)b_j.
\]

For every
\(x\notin Z\), its evaluations give a vanishing vector:
\[
 0=(f'(\xi))_{\xi\in p^{-1}(x)}
   =\sum_j f_j(x)(b_j(\xi))_{\xi\in p^{-1}(x)} .
\]
Independence on that fibre forces \(f_j(x)=0\) for
all \(j\). Thus \(f_j\in I_Z\), and the values \(b_j(\pi')\)
are integral in the common coefficient ring, so \(f'(\pi')\)
lies in \(I_Z(\pi)\). Its valuation gives
\[
 \operatorname{len}(O/I_Z(\pi))
 \le \sum_k\operatorname{len}\operatorname{Hom}_{\mathcal A}(M,H_k).
\]
If a Hom length is infinite this inequality already
holds in extended lengths. The empty generator case
has determinant one. Division by the ramification
index of \(O/\mathbb Z_\ell\) gives the normalized
inequality.  Choose a residue field
of size greater than \(g\sum_k d_k\) to specialize the
contraction simultaneously. Coefficient extraction gives
the same ideal bound over any residue field. \label{cor:congruence-determinant-2}

\end{proof}

\subsubsection*{Choosing the column cocycles}

The determinant argument requires actual columns in the integral presentation. A finite quotient specifies their evaluations only up to torsion, so we choose lifts and bound the exponent of a relation among them.

\begin{lemma}[Column lifts and their torsion tests]
\label{lem:congruence-column-cocycles}
Under the finite-class lift bound \(e'\) on the designated local
lattices, write \(c_a\) for a lift of the \(a\)-th finite quotient basis class and \(c_*\) for the extra column. Their evaluations are
\[
 c_a\longmapsto\ell^{e'}\beta_a\ (1\le a\le j_u),\qquad c_*\longmapsto\ell^{e'}\Delta_\pi\quad\text{in the unramified trivial block}.
\] These columns
satisfy the finite-class dimension bound away from \(Z\) and their
torsion dependences in the \(W\)-block are killed by \(\ell^{e'}\).
\end{lemma}

\begin{proof}
\proofstep{Integral finite bases at a column place.}
At the selected point let \(T\) be the integral lattice
in the \(W\)-block. Choose a basis \((\beta_a)_a\) of
the free quotient of \(H^1_f(E_u,T)\). This quotient
is the lattice of integral finite classes in
\(H^1_f(E_u,W)\). It spans the rational finite group
since a continuous cocycle on the compact local Galois
group has bounded image and can be scaled into \(T\).
Thus its rank is determined on the rational finite part:
\[
 \operatorname{rank}_O\bigl(H^1_f(E_u,T)/\mathrm{tors}\bigr)=\dim H^1_f(E_u,W)=j_u.
\]


Use the left ideals \(\mathcal A E_1,\mathcal A E\)
with local action by \(U_g\). Rationally these are
subrepresentations of finite direct sums of the full
twisted family, by embedding matrices as columns in
the evaluation spaces. In block \(0\), choose a
generator \(v_1\) of the image of \(E_1\) on \(T\).
Evaluation on \(v_1\) maps \(\mathcal A E_1\) onto
\(T\), because the block image is the full matrix
algebra. Evaluation of \(\mathcal A E\) on the
trivial block has image \(\Delta_\pi O\), on which
the twisted action is trivial. The two block evaluations are thus surjections of stable lattices
\[
 \mathcal A E_1\twoheadrightarrow T,\qquad \mathcal A E\twoheadrightarrow\Delta_\pi O.
\]

The twisted Galois action on the second target is trivial.
The first column lift produces finite cocycles in
\(\mathcal A E_1\) with images representing
\(\ell^{e'}\beta_a\), using lifts of the \(\beta_a\)
to integral finite classes. Use these as \(K_{u,a}\).
The second lift uses the unramified class in
\(H^1_f(E_u,\Delta_\pi O)\) with value \(\Delta_\pi\)
on arithmetic Frobenius. This class is finite since
the trivial representation's rational finite part
is the unramified line. Its lift \(K_{u,0}\) has image
with Frobenius value \(\ell^{e'}\Delta_\pi\). This
value holds on every lift of that Frobenius because
coboundaries of the trivial module vanish.

\proofstep{The local dimension at every evaluation.}
At an evaluation point \(x'\), a vector fixed by \(E_1\)
and an eigenvector \(E\eta=\Delta\eta\) give the equivariant
evaluation maps \(a\mapsto as\) on \(\mathcal A E_1\) and
\[
 a\mapsto a\eta
\]
on \(\mathcal A E\), respectively
(\(E_1s=s\)). Functoriality sends their finite cohomology
classes to finite classes. The negative Hodge jumps of the full
twisted family are constant at the column points. Removing the
cyclotomic characters at \(\pi\) removes no negative jumps.
As \(H^0(E_u,W)=0\), their common total over embeddings is
\(j_u\). The finite-part dimension formula at \(x'\) is therefore
\[
        \dim H^1_f(E_u,V)=j_u+\dim V^{D_u}.
\]
This is the local dimension used for the finite column cocycles.

\proofstep{Torsion relations in the chosen basis.}
The rational dimension leaves the exponent of integral torsion relations undetermined.

Test a torsion relation on the \(W\)-block. Write a linear combination
of the evaluated column cocycles with coefficients
in \(F/O\), and lift those coefficients to \(F\).
If the combination is a coboundary modulo \(T\),
the rational finite class \(\sum_a y_a\ell^{e'}\beta_a\)
comes from \(H^1(E_u,T)\) by exactness for
\(0\to T\to W\to W/T\to0\).
It therefore comes from \(H^1_f(E_u,T)\) by the
definition of the integral finite part. Because
the \(\beta_a\) form its free basis, each
\(y_a\ell^{e'}\) belongs to \(O\). This proves the
killing exponent for such a dependence.   \label{cor:congruence-determinant-3}

\end{proof}

\subsubsection*{The Selmer length comparison for this presentation}
\label{subsec:congruence-selmer}

We can now evaluate the presentation at the chosen theta point. Each block contributes a Selmer length, while the valuation of the specialized determinant records the cost of imposing the integral conditions.

At the chosen point let \(\Delta_\pi\ne0\) be the
specialization of \(\Delta\) and put \(v_\Delta
=v_\ell(\Delta_\pi)\). From now on lengths are
normalized by \(v_\ell(\ell)=1\).
The block targets are \(H_0=(F/O)^r\) with the
action of \(W\), \(H_1=F/O\) with trivial action,
and \(H_k=(F/O)(1-k)\) for \(2\le k\le t\).
Write \(P_k=\operatorname{Hom}_{\mathcal A}(M,H_k)\).
\begin{proposition}[From block lengths to Selmer kernels]
\label{prop:congruence-block-selmer-length}
With the integral presentation and column lifts of
Lemma~\ref{lem:congruence-column-cocycles}, the normalized sum of
block Hom lengths in Proposition~\ref{prop:theta-determinant-length}
is bounded by the global kernel lengths and the local and
presentation costs below. The three global kernels are
the rank-\(r\) finite kernel, the unramified first character and
the strict classes in twists \(1-k\) for \(2\le k\le t\).
If \(c_k\) is the normalized contribution charged to coboundaries in block \(k\), then
\[
 c_1=0,\qquad c_k\le2d_kv_\Delta\quad(k\ne1).
\]
\end{proposition}

\begin{proof}
\proofstep{The map to continuous cocycles.}
A homomorphism \(\varphi\) supplies the cocycle
\[
    c_\varphi(g)=\varphi((U_g-1)E).
\]
The cocycle identity follows by multiplying \(U_g U_h\).
In the integral presentation all the matrix functions
are continuous, and the orders and their modules are
finite over \(O\). A map to \(H_k\) kills a common
power of the uniformizer times the module because the
images of a finite set of generators are torsion.
Thus \(c_\varphi\) is continuous.

\proofstep{Its kernel on the local generators.}
First consider the map from \(P_k\) to cocycles before
taking cohomology. Its kernel consists of maps vanishing
on \(Y\). They are determined by the images of the local
generators, subject to the relations of \(M\).
The specializations on the target blocks satisfy
\[
 E|_{H_0}=0,\qquad E|_{H_1}=\Delta_\pi,\qquad E|_{H_k}=0\quad(k\ge2).
\]
The idempotent \(E_1\) has rank one
in block \(0\) and is zero on the character blocks.
The extra column cocycle \(K_{u,0}(g)\) lies in the
left ideal generated by \(E\).

\begin{itemize}
\item In block \(0\) the column lift input says that
a linear combination \(\sum_a K_{u,a}(g)s_a\) with
\(s_a\in E_1H_0\) which is a coboundary on \(D_u\)
forces \(\ell^{e'}s_a=0\) for every \(a\).
For a map vanishing on \(Y\) the combination with
coefficients \(\Delta_\pi\varphi(s_{u,a})\) is a
coboundary by the relation for \(D_u\). The tuple of images of the \(s_{u,a}\) lies in
\[
   (E_1 H_0)[\ell^{e'}\Delta_\pi]^{\,j_u},
\]
of normalized length at most \(j_u(e'+v_\Delta)\).

For a fixed tuple, differences of choices of \(n_u\) lie in
\(H_0^{D_u}\), of length \(h_W(u)\), and the remaining
image of \(s'_{u,0}\) lies in \(H_0[\Delta_\pi]\), of length
at most \(rv_\Delta\). Kernel plus image for projection onto
the tuple therefore gives
\(j_u(e'+v_\Delta)+h_W(u)+rv_\Delta\).

\item In block \(1\) the matrix \(K_{u,0}\) has some
value on \(D_u\) acting by
\(\ell^{e'}\Delta_\pi\) times a unit, the first-character
column input. The relation for that element gives
\[
 \ell^{e'}\Delta_\pi^2\varphi(s_{u,0}')=0
\]
for a
map vanishing on \(Y\). The image of \(n_u\) is killed
by \(\Delta_\pi\) and that of \(s_{u,a}\) is zero.
The resulting cost is at most \(e'+3v_\Delta\).

\item In block \(k\ge2\) the \(s_{u,a}\) vanish and
\(\Delta_\pi\varphi(s_{u,0}')=0\). The relation for
\(D_u\) forces \(\varphi(n_u)\in H_k^{D_u}\), since
\(K_{u,0}\) acts by zero. With the scalar invariant
bound \(h_0\) the cost is \(h_0+v_\Delta\).
\end{itemize}

\proofstep{The cost of quotienting by coboundaries.}
Coboundaries change the length differently on each block.

Pass next to cohomology. In block \(1\) the coboundaries
are zero. In the other blocks the central scalars \(a_g\) of the marked corner satisfy
\[ E=\sum_g a_g U_g,\qquad \sum_g a_g=\Delta_\pi\quad\text{at the chosen point}. \]
If \(c_\varphi(g)=(U_g-1)v\), the element
\[
    y=\sum_g a_g (U_g-1)E=(E-\textstyle\sum_g a_g)E
\]
has image \(-\Delta_\pi v\) in \(H_k\), since \(E\)
acts by zero on \(H_k\). The algebra identity
\(Ey=\Delta y\) then forces \(\Delta_\pi^2 v=0\).
The space of image cocycles that are coboundaries
has length at most \(2d_k v_\Delta\).

\proofstep{Relative Selmer kernels and local correction terms.}
Finally let \(C_k\) be the image of \(P_k\) in
\(H^1(E,H_k)\). At each local place impose a target
subgroup \(F_{k,u}\) of local cohomology. Suppose
the restrictions of \(C_k\) are in subgroups
\(B_{k,u}\supset F_{k,u}\) and that the classes also
lie in a global subgroup \(Q_k\). The kernel of
\[
   C_k\longrightarrow\bigoplus_u B_{k,u}/F_{k,u}
\]
lies in the subgroup of \(Q_k\) satisfying all the
local conditions. Thus its global length bound and
the sum of the local colength bounds give a length
bound for \(C_k\). In the split comparison, the
column images already satisfy \(F_{k,u}\).

In
block \(0\) the row quotient has the colength
bound proved above, and the global kernel has length
at most \(s_\pi\). In the character blocks the row and
tame quotients have bounds \(([E_w:\mathbb Q_\ell]+2)c_B\)
and \(v_\ell(\zeta_v-1)\). If \(K_1\) is the global kernel in block one, unramified class field theory gives
\[
 K_1\hookrightarrow\operatorname{Hom}(\operatorname{Cl}(E),F/O),\qquad \operatorname{len}_{\rm norm} K_1\le v_\ell(h_E).
\]
In the higher blocks the local
conditions defining this kernel are strict at \(\ell\)
and unramified elsewhere. Its length is bounded by the
strict group \(\operatorname{Sha}^1(E,(F/O)(1-k))\).
The normalization and the row and tame restriction arguments
follow below.

\proofstep{The sum of the normalized block costs.}
Let \(k_c\) be the number of columns. Summing these
bounds over all blocks gives the expression used in
the split-row estimate. The terms from the rank-\(r\)
and character blocks contribute
\[
 \sum_u\{j_u(e'+v_\Delta)+h_W(u)+e'\},
\]
and the \(t-1\) higher blocks contribute at most
\(tk_ch_0\). The remaining denominator terms from
the presentation and the coboundaries are bounded by
\(2k_cv_\Delta+4R(k_c+1)v_\Delta\). Indeed the dimensions \(d_k\) of the blocks sum to \(R\), so their contribution is
\[
 \sum_k d_k k_c v_\Delta=R k_c v_\Delta,
\]
and the coboundaries cost at most \(2Rv_\Delta\).

The row and tame costs of the
character blocks occur at most \(t\) times. Adding the
global kernels accounts for \(s_\pi\), the class
number and the strict cohomology sum. Pairing the row
cost for \(W\) at \(w\) with its conjugate column
\(u\) and using \(j_u\le [E_u:\mathbb Q_\ell]m\)
gives the formula below.
\label{cor:congruence-determinant-4}
\end{proof}
